% ============================================================
% SECTION 09
% SCOPE AND LIMITS OF INVARIANCE
% Chapter 02: Invariance
% Book I: Ontology of Viability
% ============================================================

\section{Scope and Limits of Invariance}
\label{sec:scope-and-limits-of-invariance}

% ------------------------------------------------------------
% 9.1. THE NEED TO ESTABLISH THE SCOPE OF INVARIANCE
% ------------------------------------------------------------

The preceding sections have developed a provisional
relational account of invariance and examined its
relationship with identity, difference, change,
continuity, structure, different forms of determination,
and knowledge.

These examinations have clarified the analytical
use of invariance, but they have also established
the need to delimit its scope.

A concept cannot be used adequately if the conditions
under which it applies remain unspecified. Invariance
is particularly vulnerable to this problem because
claims of invariance may be extended beyond the
determinations, transformations, and domains for
which they have been established.

The statement that something is invariant is never
sufficient by itself. It must be possible to identify
what is invariant, under which transformations,
and in what respect it remains unchanged.

The present section consolidates these requirements
and distinguishes the commitments adopted in this
chapter from the questions that remain unresolved.

Its purpose is not to establish an exhaustive
ontology of invariance. Rather, it clarifies the
legitimate scope of the concept within the
provisional framework of the first edition.

% ------------------------------------------------------------
% 9.2. THE PROVISIONAL RELATIONAL DEFINITION
% ------------------------------------------------------------

For the purposes of this first edition, invariance
is defined as a relation between an existent, or
a determinate aspect of an existent, and a specified
domain of changes or transformations, such that
the existent or aspect remains unchanged in a
particular respect under those changes or
transformations.

This definition establishes three analytical
components:

\begin{enumerate}
    \item \textbf{The subject of invariance:}
    the existent or determinate aspect whose
    invariance is being examined.

    \item \textbf{The domain of transformation:}
    the specified changes or transformations
    in relation to which invariance is asserted.

    \item \textbf{The respect of invariance:}
    the particular determination, feature,
    or relation that remains unchanged.
\end{enumerate}

These components are analytical distinctions
within a single invariance claim. They must
not be interpreted as three independently
existing ontological entities.

The definition is intentionally relational.
It provides a framework for identifying and
examining invariance without presupposing
that invariance is itself an independent
substance, an immutable foundation, or
a universal property of every existent.

The relational definition is a methodological
and conceptual commitment of the present
edition. It does not constitute a final
metaphysical conclusion concerning the
ultimate ontological status of invariance.

% ------------------------------------------------------------
% 9.3. THE REQUIREMENT OF SPECIFICATION
% ------------------------------------------------------------

Every claim of invariance must be sufficiently
specified to distinguish it from an unrestricted
assertion of permanence or immutability.

The subject, domain, and respect of invariance
must be identifiable, either explicitly or
through the context in which the claim is made.

An assertion that a structure remains invariant,
for example, is incomplete unless it is possible
to determine which structural aspect is being
considered and which transformations are
included in the relevant domain.

The same requirement applies to quantitative,
qualitative, relational, and other forms of
invariance.

The scope of an invariance claim must also
remain proportionate to the grounds on which
it is established.

The recognition of invariance within a limited
domain does not justify its extension to
transformations that have not been considered.

Accordingly, the following inference is not
generally justified:

\[
\text{Invariance within } \mathcal{D}_1
\not\Rightarrow
\text{Invariance within every domain}
\]

The expression indicates that the identification
of invariance within one specified domain does
not establish invariance under all possible
transformations.

This limitation is not a weakness of the concept.
It is a condition of its precise application.

% ------------------------------------------------------------
% 9.4. INVARIANCE IS NOT IDENTITY
% ------------------------------------------------------------

The distinction between invariance and identity
is one of the central commitments of this chapter.

Identity concerns the identity-bearing character
of an existent, which, according to the foundational
position of this book, is ontologically inseparable
from being.

Invariance, as provisionally defined here,
concerns the persistence of a specified
determination under a specified domain of
transformations.

The two concepts therefore address different
questions.

The identity of an existent cannot be reduced
to the invariance of one or more of its
determinations.

Likewise, the identification of an invariant
determination does not, by itself, establish
the identity of the existent.

The distinction also prevents a premature
conclusion that every transformation of an
existent must either preserve or destroy
its identity according to a single invariant
criterion.

The relationship between identity and
transformation requires a broader ontological
account than the present chapter provides.

The origin and deeper constitution of identity
remain deferred to the subsequent development
of the realms of being, light, and ultimately
the dynamics of being.

% ------------------------------------------------------------
% 9.5. INVARIANCE IS NOT IMMUTABILITY
% ------------------------------------------------------------

Invariance must also be distinguished from
immutability.

Immutability, in its unrestricted sense,
suggests the absence of change. The relational
concept of invariance developed here does
not require the absence of every possible
change.

An existent may undergo transformations
while a particular determination remains
unchanged within the specified domain.

Consequently, the identification of an
invariant aspect does not establish that
the existent is immutable as a whole.

Nor does the absence of invariance in
one respect establish that the existent
is subject to unrestricted change in
every other respect.

The distinction prevents the concept
of invariance from being used to introduce
an unexamined metaphysics of absolute
permanence.

It also prevents change from being treated
as the universal negation of invariance.

The relevant question remains whether
a particular determination is unchanged
under the transformations being considered.

% ------------------------------------------------------------
% 9.6. INVARIANCE IS NOT CONTINUITY
% ------------------------------------------------------------

The distinction between invariance and
continuity must remain equally explicit.

Invariance concerns the persistence of
a specified determination under a defined
domain of transformations.

Continuity concerns whether and how an
existent continues through different
conditions or phases.

The invariance of a particular feature
may be relevant to an examination of
continuity, but it does not independently
establish that continuity.

Conversely, the continuation of an
existent does not necessarily require
the unchanged persistence of every
determination associated with it.

The concept of continuity therefore
cannot be reduced to the persistence
of invariant features.

This distinction also prevents the
present account from presupposing
a permanent substrate or fixed essence
as the necessary ground of continuation.

The ontological relationship between
identity, continuity, and transformation
remains open for examination within
the broader development of the system.

% ------------------------------------------------------------
% 9.7. INVARIANCE AND STRUCTURE:
% THE LIMIT OF GENERALIZATION
% ------------------------------------------------------------

The analysis of structural invariance
has established that the persistence
of a particular structural determination
does not necessarily entail the invariance
of the entire structure.

A relation may remain unchanged while
other relations or determinations are
altered.

Likewise, a structure may preserve
a specified organization under some
transformations without remaining
unchanged in every respect.

The distinction between local and
general invariance is therefore essential.

The invariance of a limited aspect
must not be generalized to the whole
existent or its entire organization
without adequate grounds.

This limitation applies not only to
structural invariance but also to
quantitative, qualitative, and relational
invariance.

No particular form of invariance
automatically establishes a universal
claim concerning the existent as a whole.

The scope of every claim must remain
tied to the determination and domain
that have actually been examined.

% ------------------------------------------------------------
% 9.8. EPISTEMIC LIMITS AND ONTOLOGICAL STATUS
% ------------------------------------------------------------

The identification of invariance is
epistemically mediated.

A knower must identify the relevant
determination, specify the domain of
transformations, and establish the
grounds for claiming that something
remains unchanged.

These operations may involve observation,
comparison, measurement, formalization,
or conceptual analysis.

However, the epistemic conditions
through which invariance is recognized
must not be confused with its ontological
status.

The absence of knowledge concerning
a particular invariant does not establish
that no such invariance exists.

Similarly, the recognition of invariance
within one domain does not establish
its presence within every other domain.

Knowledge and non-knowledge are
intertwined within the situated condition
of the knower. Knowledge of one aspect
of an existent may coexist with
non-knowledge of other aspects that
could be relevant to the interpretation
of its invariance.

This situation does not justify either
an unrestricted skepticism or an
unqualified claim of knowledge.

It requires that ontological assertions
and epistemic justifications be
distinguished and appropriately
related.

The existent is not reduced to what
a particular knower recognizes, but
claims concerning the existent remain
subject to the epistemic grounds
through which they are justified.

% ------------------------------------------------------------
% 9.9. THE OPEN ONTOLOGICAL STATUS OF INVARIANCE
% ------------------------------------------------------------

The relational definition adopted in
this first edition provides a workable
framework for examining invariance
across different determinations and
domains of transformation.

Nevertheless, the question remains
whether invariance possesses a deeper
ontological significance that cannot
be fully expressed through its
relational formulation.

The present chapter does not resolve
this question.

It would be premature to conclude
that invariance is merely an epistemic
construct, because the relational
definition concerns an existent or
one of its determinations and does
not reduce its status to the act
of knowing.

It would be equally premature to
conclude that invariance constitutes
an independent ontological principle
or a fundamental condition of being.

Neither conclusion has been
established by the analysis developed
so far.

The decision to employ a relational
definition in the first edition
is therefore a provisional
methodological commitment.

Its adequacy may be reconsidered
in subsequent revisions, when
the broader ontological framework
has been developed sufficiently
to examine whether invariance
has a deeper status.

This openness is deliberate.
It preserves the distinction between
what the present framework can
legitimately establish and what
remains a question for further
ontological investigation.

% ------------------------------------------------------------
% 9.10. ESTABLISHED COMMITMENTS AND
% UNRESOLVED QUESTIONS
% ------------------------------------------------------------

The analysis of this chapter
has established a number of
conceptual and methodological
commitments.

These commitments can be
distinguished from questions
that remain unresolved.

\subsection*{Established Commitments}

The following principles govern
the use of invariance within
the present chapter:

\begin{enumerate}
    \item Invariance is provisionally
    defined as a relation between
    an existent or a determinate
    aspect of an existent and
    a specified domain of
    transformations.

    \item Every invariance claim
    requires the specification
    of its subject, domain,
    and respect.

    \item Invariance is not
    identical to identity.

    \item Invariance is not
    equivalent to immutability
    or the absence of all change.

    \item Invariance is not
    identical to continuity.

    \item The invariance of
    a particular determination
    does not automatically
    establish the invariance
    of the whole existent.

    \item Different forms and
    domains of invariance may
    overlap and are not
    necessarily mutually
    exclusive.

    \item The epistemic recognition
    of invariance must be
    distinguished from its
    ontological status.

    \item Knowledge and
    non-knowledge must both
    be considered when the
    epistemic justification
    of invariance claims
    is examined.

    \item The relational
    definition is a provisional
    commitment of the first
    edition rather than a
    final metaphysical
    conclusion.
\end{enumerate}

\subsection*{Unresolved Questions}

The following questions remain
open and must not be treated
as settled conclusions:

\begin{enumerate}
    \item Does invariance possess
    an ontological significance
    beyond its relational
    formulation?

    \item What is the deeper
    relationship between
    invariance and the
    identity-bearing character
    of existence?

    \item Can invariance be
    understood as a condition
    of certain forms of
    continuity, or is its
    relevance always dependent
    upon the specific domain
    being examined?

    \item How does invariance
    relate to the different
    realms of being?

    \item Does the ontological
    status of invariance
    differ across different
    realms or forms of
    manifestation?

    \item What is the relationship
    between invariance and
    the dynamics of being?

    \item Is there a deeper
    ground that explains
    why certain determinations
    remain invariant through
    particular transformations?
\end{enumerate}

These questions identify the
limits of the present account.
They are not reasons to abandon
the relational definition,
but reasons to preserve its
provisional status.

% ------------------------------------------------------------
% 9.10.1. TABLE: SCOPE AND LIMITS OF INVARIANCE
% ------------------------------------------------------------


% ============================================================
% TABLE 09 — SCOPE AND LIMITS OF INVARIANCE
% Book I: Ontology of Viability
% Part I: Identity and Manifestation
% Chapter 2: Invariance
% ============================================================

\begingroup

\small
\renewcommand{\arraystretch}{1.25}

\setlength{\LTleft}{0pt}
\setlength{\LTright}{0pt}
\setlength{\LTpre}{1.0em}
\setlength{\LTpost}{1.0em}

\begin{longtable}{
    @{}
    >{\raggedright\arraybackslash}p{0.18\textwidth}
    >{\raggedright\arraybackslash}p{0.36\textwidth}
    >{\raggedright\arraybackslash}p{0.36\textwidth}
    @{}
}

\caption{The Scope and Limits of the Concept of Invariance}
\label{tab:scope-and-limits-of-invariance}
\\

\toprule
\textbf{Dimension} &
\textbf{What the Concept Establishes} &
\textbf{What It Does Not Establish}
\\
\midrule
\endfirsthead

\multicolumn{3}{@{}l@{}}{
    \textit{The Scope and Limits of the Concept of Invariance (continued)}
}
\\
\toprule
\textbf{Dimension} &
\textbf{What the Concept Establishes} &
\textbf{What It Does Not Establish}
\\
\midrule
\endhead

\midrule
\multicolumn{3}{r@{}}{
    \textit{Continued on the next page}
}
\\
\endfoot

\bottomrule
\endlastfoot

Subject of invariance &
An existent, or a determinate aspect of an existent,
may be considered with respect to invariance. &
The subject of invariance is not thereby reduced
to the particular determination under consideration.
\\

\addlinespace[0.5em]

Domain of change &
Invariance is assessed in relation to a specified
domain of changes or transformations. &
Invariance within one domain does not establish
invariance under every possible change
or transformation.
\\

\addlinespace[0.5em]

Respect of invariance &
A particular determination, feature, or relation
may remain unchanged under the specified
transformations. &
The invariance of one respect does not establish
the invariance of every aspect of the existent.
\\

\addlinespace[0.5em]

Identity &
Invariance may help describe how a determinate
aspect remains unchanged under specified
transformations. &
Invariance does not constitute identity, provide
its complete explanation, or establish that
identity is exhausted by invariant features.
\\

\addlinespace[0.5em]

Difference &
Different determinations may coexist with
the invariance of a particular determination. &
Difference does not necessarily negate invariance,
and invariance does not eliminate difference.
\\

\addlinespace[0.5em]

Change &
An existent may undergo change in some respects
while remaining invariant in another respect. &
Invariance does not mean the total absence
of change or the immutability of the existent
as a whole.
\\

\addlinespace[0.5em]

Continuity &
Invariance may contribute to the description
of certain patterns of continuity. &
The invariance of a particular feature does not,
by itself, establish the continuity of the existent
as a whole.
\\

\addlinespace[0.5em]

Structure &
Particular structural determinations or relations
may remain invariant under specified
transformations. &
The invariance of one structural relation does not
establish the invariance of the entire structure.
\\

\addlinespace[0.5em]

Knowledge &
Invariance may be investigated, recognized,
compared, and formally described through
epistemic operations. &
Epistemic recognition does not constitute
ontological invariance, and the absence
of recognition does not establish its absence.
\\

\addlinespace[0.5em]

Ontological significance &
The relational concept provides an initial
framework for describing determinate invariance. &
This framework does not settle whether invariance
has a deeper ontological status beyond the
relational account adopted in this edition.
\\

\end{longtable}

\endgroup

% ============================================================
% END OF TABLE 09
% ============================================================


% ------------------------------------------------------------
% 9.11. SYNTHESIS AND TRANSITION
% ------------------------------------------------------------

The present section has clarified
the scope within which invariance
can be used in the ontology of
viability.

The relational definition provides
a precise analytical framework
for examining what remains
unchanged under specified
transformations.

Its use requires the identification
of the existent or determination
under consideration, the domain
of transformation, and the
particular respect in which
invariance is asserted.

The analysis has also established
that invariance cannot be reduced
to identity, immutability,
continuity, structure, or
epistemic recognition.

These concepts are related,
but each addresses a distinct
problem.

The limits identified throughout
the chapter are essential to
preventing the concept of
invariance from carrying
metaphysical commitments
that have not yet been justified.

The relational definition is
therefore retained as the
working framework of the
first edition, while its
possible deeper ontological
significance remains open
for subsequent reconsideration.

The final section of this chapter
will synthesize the conceptual
results, consolidate the
established distinctions,
and identify the role of
invariance within the broader
development of the ontology
of viability.

% ============================================================
% END OF SECTION 09
% ============================================================